% !TeX program = xelatex
% Evensong 研究提案 v2 — EverMind 研究资助申请
% Compile: latexmk -xelatex evensong-research-proposal-v2-zh.tex
% Version: 2.0 (2026-04-10) — 新增方法验证协议、自修复分类、事故记录、MiniMax-M2.7

\documentclass[11pt, a4paper]{article}

% ═══════════════════════════════════════════════════════════
% § TYPOGRAPHY & LAYOUT — The Invisible 80%
% ═══════════════════════════════════════════════════════════
\usepackage[
  top=25mm, bottom=28mm,
  inner=28mm, outer=32mm,
  headheight=14pt, headsep=10mm
]{geometry}

\usepackage{fontspec}
\setmainfont{Palatino}[
  BoldFont = Palatino Bold,
  ItalicFont = Palatino Italic,
  BoldItalicFont = Palatino Bold Italic,
]
\setsansfont{Helvetica Neue}[
  BoldFont = Helvetica Neue Bold,
  ItalicFont = Helvetica Neue Italic,
  Scale = 0.95,
]
\setmonofont{Menlo}[Scale=0.82]

\usepackage{xeCJK}
\setCJKmainfont{Songti SC}[BoldFont=Songti SC Bold]
\setCJKsansfont{PingFang SC}

\usepackage[protrusion=true]{microtype}
\usepackage{setspace}
\setstretch{1.35}
\setlength{\parindent}{0pt}
\setlength{\parskip}{0.6em}

% ═══════════════════════════════════════════════════════════
% § COLOR SYSTEM — DASH SHATTER brand palette
% ═══════════════════════════════════════════════════════════
\usepackage[dvipsnames,table,svgnames]{xcolor}
% Primary palette from DASH SHATTER theme.ts
\definecolor{deep}{RGB}{53,88,76}         % #35584C — forest green, primary brand
\definecolor{accent}{RGB}{240,238,155}    % #F0EE9B — neon yellow, highlights
\definecolor{warm}{RGB}{255,193,7}        % #FFC107 — warning yellow
\definecolor{success}{RGB}{86,190,120}    % #56BE78 — positive metrics
\definecolor{subtle}{RGB}{238,235,225}    % #EEEBE1 — warm cream wash
\definecolor{faded}{RGB}{127,136,130}     % #7F8882 — secondary text
\definecolor{cream}{RGB}{248,245,239}     % #F8F5EF — page background
\definecolor{deepdark}{RGB}{16,41,31}     % #10291F — deep forest for Investment box
\definecolor{glass}{RGB}{53,88,76}        % glassmorphism base (used at low opacity)
\definecolor{errcoral}{RGB}{255,107,128}  % #FF6B80 — error/danger accent

% ═══════════════════════════════════════════════════════════
% § SECTION STYLING — titlesec
% ═══════════════════════════════════════════════════════════
\usepackage{titlesec}

\titleformat{\section}
  {\sffamily\Large\bfseries\color{deep}}
  {\thesection}{0.7em}{}
  [\vspace{-0.4em}{\color{accent}\hrule height 0.8pt}\vspace{0.15em}]

\titleformat{\subsection}
  {\sffamily\large\color{deep}}
  {\thesubsection}{0.6em}{}

\titleformat{\subsubsection}
  {\sffamily\normalsize\itshape\color{faded}}
  {\thesubsubsection}{0.5em}{}

\titlespacing*{\section}{0pt}{2ex plus 0.5ex}{1ex}
\titlespacing*{\subsection}{0pt}{1.5ex plus 0.3ex}{0.6ex}

% ═══════════════════════════════════════════════════════════
% § HEADERS & FOOTERS
% ═══════════════════════════════════════════════════════════
\usepackage{fancyhdr}
\pagestyle{fancy}
\fancyhf{}
\fancyhead[L]{\small\sffamily\color{faded}\itshape Evensong Benchmark Framework}
\fancyhead[R]{\small\sffamily\color{faded}Research Proposal — EverMind}
\fancyfoot[C]{\small\sffamily\color{faded}\thepage}
\renewcommand{\headrulewidth}{0.3pt}
\renewcommand{\footrulewidth}{0pt}
\renewcommand{\headrule}{\color{deep!40}\hrule height 0.3pt}

% ═══════════════════════════════════════════════════════════
% § CALLOUT BOXES — tcolorbox
% ═══════════════════════════════════════════════════════════
\usepackage[most]{tcolorbox}

% ── Glassmorphism: deep green glass with high-contrast titles ──
% Key Finding box — deep green title bar, WHITE title text (high contrast)
\newtcolorbox{keyfinding}[1][]{
  enhanced, breakable,
  colback=deep!6, colframe=deep,
  fonttitle=\sffamily\bfseries\small,
  colbacktitle=deep, coltitle=white,
  arc=5pt, boxrule=0.8pt,
  shadow={1pt}{-1pt}{0pt}{deep!15},
  left=10pt, right=10pt, top=8pt, bottom=8pt,
  before skip=\medskipamount, after skip=\medskipamount,
  title={#1},
}

% Investment Thesis box — deep forest body, deep green title bar with WHITE text
\newtcolorbox{investmentthesis}[1][]{
  enhanced, breakable,
  colback=deepdark, colframe=deep,
  fonttitle=\sffamily\bfseries\Large,
  colbacktitle=deep, coltitle=white, coltext=white,
  arc=6pt, boxrule=1.2pt,
  shadow={2pt}{-2pt}{0pt}{deepdark!40},
  left=14pt, right=14pt, top=12pt, bottom=12pt,
  before skip=\bigskipamount, after skip=\bigskipamount,
  title={#1},
}

% Evidence box — light green tint body, deep green title bar, WHITE title
\newtcolorbox{evidence}[1][]{
  enhanced, breakable,
  colback=deep!5, colframe=deep!50,
  fonttitle=\sffamily\bfseries\small,
  colbacktitle=deep!80, coltitle=white,
  arc=4pt, boxrule=0.5pt,
  shadow={1pt}{-1pt}{0pt}{deep!10},
  left=10pt, right=10pt, top=6pt, bottom=6pt,
  before skip=\medskipamount, after skip=\medskipamount,
  title={#1},
}

% Metric highlight (inline)
\newtcolorbox{metric}{
  enhanced, on line, arc=3pt, boxrule=0pt,
  colback=accent!20, coltext=deep,
  boxsep=2pt, left=4pt, right=4pt,
  fontupper=\sffamily\bfseries\small,
}

% ═══════════════════════════════════════════════════════════
% § PreTeXt CONTAINMENT PATTERN — adjustbox + linewidth sizing
% Adopted from PreTeXt's LaTeX output methodology:
% all block content constrained to \linewidth, no overflow
% ═══════════════════════════════════════════════════════════
\usepackage[export]{adjustbox}  % max width, center, valign for overflow prevention

% ═══════════════════════════════════════════════════════════
% § TABLES — tabularray + siunitx
% ═══════════════════════════════════════════════════════════
\usepackage{tabularray}
\UseTblrLibrary{booktabs}
\usepackage{siunitx}
\sisetup{group-separator={,}, group-minimum-digits=4}

% ═══════════════════════════════════════════════════════════
% § FIGURES & CHARTS
% ═══════════════════════════════════════════════════════════
\usepackage{pgfplots}
\pgfplotsset{compat=1.18}
\usepgfplotslibrary{fillbetween}
\usepackage{tikz}
\usetikzlibrary{positioning, arrows.meta, shapes.geometric, fit, calc, backgrounds}
\usepackage[font=small, labelfont=bf, textfont=it,
            labelsep=period, justification=raggedright]{caption}
\usepackage{subcaption}
\usepackage{float}

% ═══════════════════════════════════════════════════════════
% § UTILITIES
% ═══════════════════════════════════════════════════════════
\usepackage{enumitem}
\setlist{nosep, leftmargin=1.5em}
\usepackage{hyperref}
\hypersetup{
  colorlinks=true,
  linkcolor=deep,
  urlcolor=deep,
  citecolor=deep,
  pdftitle={Evensong: Memory-Augmented AI Agent Benchmark Framework},
  pdfauthor={Nolan Zhu, DASH SHATTER Research},
}
\usepackage{bookmark}
\usepackage{graphicx}
\usepackage{amsmath,amssymb}

% Custom commands
\newcommand{\runid}[1]{\texttt{\color{deep}#1}}
\newcommand{\numval}[1]{\textbf{\color{success!70!black}#1}}
\newcommand{\metricval}[1]{{\sffamily\bfseries\color{success!70!black}#1}}

% ═══════════════════════════════════════════════════════════
%                     DOCUMENT BEGIN
% ═══════════════════════════════════════════════════════════
\begin{document}
\thispagestyle{empty}

% ── TITLE PAGE ─────────────────────────────────────────────
\begin{center}
\vspace*{35mm}

{\color{deep}\rule{0.6\textwidth}{1pt}}

\vspace{8mm}

{\fontsize{28pt}{34pt}\selectfont\sffamily\bfseries\color{deep}%
Evensong}

\vspace{4mm}

{\fontsize{13pt}{18pt}\selectfont\sffamily\color{faded}%
记忆增强评估框架\\[2pt]
面向自主 AI 编程智能体}

\vspace{8mm}

{\color{deep}\rule{0.6\textwidth}{1pt}}

\vspace{14mm}

{\large\sffamily\color{deep!80} 研究提案}

\vspace{3mm}

{\sffamily\color{faded} 提交至 \textbf{EverMind} 申请研究资助}

\vspace{20mm}

\begin{adjustbox}{max width=\linewidth}
\begin{tblr}{
  colspec = {rl},
  column{1} = {font=\sffamily\small\color{faded}},
  column{2} = {font=\sffamily},
  rowsep = 3pt,
}
Principal Investigator & Nolan Zhu (朱恒源) \\
Affiliation & DASH SHATTER Research \\
Infrastructure Partner & EverMind / EverOS \\
Date & April 2026 \\
Classification & Confidential \\
\end{tblr}
\end{adjustbox}

\vfill

{\small\sffamily\color{faded}%
\textit{``AI agent memory causally changes engineering decisions,}\\
\textit{and pressure is required to trigger self-evolution.''}\\[4pt]
--- Empirical finding from Evensong R011-B (2026-04-10)}

\vspace{10mm}
\end{center}

\clearpage

% ── EXECUTIVE SUMMARY ──────────────────────────────────────
\pagestyle{fancy}
\section*{执行摘要}
\addcontentsline{toc}{section}{执行摘要}

\begin{keyfinding}[核心发现]
经过 13 轮基准测试 (R001--R012, R006-Grok)，逾 6{,}000 个自动化测试，我们首次提供了\textbf{持久记忆因果性地改变 AI 智能体工程策略}的实证证据，并发现\textbf{环境压力是自主自我进化的必要条件}。上述发现依托 EverOS 的记忆基础设施实现，对下一代 AI 智能体架构具有直接启示。
\end{keyfinding}

\textbf{Evensong} 框架是一个自动化基准测试工具，用于在受控记忆和压力条件下评估 AI 编码智能体。基于 DASH SHATTER（逆向工程、完全可 hack 的 Claude Code CLI）构建，测量智能体在微服务架构生成、测试质量和涌现行为进化方面的表现。

\textbf{关键成果：}
\begin{itemize}
  \item \textbf{460\% 测试产出增长}：从 R005 到 R009（$80 \to 1{,}051$ 个测试），驱动力为迭代式提示工程进化
  \item \textbf{首次记录递归记忆污染}：EverOS 检索将策略反馈给智能体，智能体再将实验知识回写至记忆，形成自放大反馈环路
  \item \textbf{2\texorpdfstring{$\times$}{x}2 因子实验设计}：隔离记忆$\times$压力交互效应 --- 这是首个针对 AI 智能体自我进化的受控实验
  \item \textbf{9 模型基准矩阵}：覆盖 Claude、GPT、Grok、Gemini、GLM、Qwen、DeepSeek、Kimi 和 MiniMax-M2.7
  \item \textbf{新颖的跨模型协作维度}：现有基准测试均未涉及 AI 与 AI 之间的协调能力评估
  \item \textbf{4 型自修复分类法}：首个对 AI 智能体错误自我纠正策略的分类体系
  \item \textbf{60 倍成本验证协议}：用便宜模型验证方法可防止 80 美元以上的浪费
\end{itemize}

\textbf{资金申请：}我们提议与 EverMind 继续合作，以 EverOS 作为 Evensong 的记忆骨干，目标是建立首个记忆增强型 AI 智能体开放基准测试，并在 SWE-bench Multimodal 和 HAL Agent Leaderboard 上进入前三名。

\clearpage
\tableofcontents
\clearpage

% ══════════════════════════════════════════════════════════════
% SECTION 1: BACKGROUND & MOTIVATION
% ══════════════════════════════════════════════════════════════
\section{背景与动机}

\subsection{AI 智能体评估中缺失的变量}

现有 AI 编码基准测试（SWE-bench、HumanEval、MBPP）将智能体评估为无状态函数：给定提示词，产出代码。这忽略了一个关键维度——真实世界智能体部署中的\textbf{持久记忆}。

生产级 AI 智能体跨会话积累战略性知识：哪些架构模式成功、哪些调试方法失败、哪些测试结构能捕获边界情况。这种积累的上下文从根本上改变了智能体处理工程问题的方式。然而，目前没有任何基准测试测量这种效应。

\subsection{EverOS 作为记忆基础设施}

EverMind 的 EverOS 为研究记忆对 AI 智能体的影响提供了理想基础设施：

\begin{itemize}
  \item \textbf{记忆空间}：物理隔离的存储域（我们使用 3 个：\texttt{zonicdesign.art}、\texttt{allaround}、\texttt{void}），实现受控记忆条件
  \item \textbf{双模型提取}：边界检测（GPT-4.1-mini）+ 语义提取（Qwen3-235B）管道将原始对话转换为结构化 MemCell
  \item \textbf{代理检索}：多轮查询扩展，呈现策略级记忆，而非仅关键词匹配
  \item \textbf{发送方身份}：区分人类观察者、AI 观察者和 AI 运行者——对于追踪记忆污染路径至关重要
  \item \textbf{多模态存储}：截图和视觉证据与文本记忆一起存储
  \item \textbf{Cases \& Skills}：形成的可复用智能体能力的提取模式
\end{itemize}

\begin{evidence}[为什么选 EverOS 而非自建记忆系统？]
我们最初考虑为基准测试记忆隔离构建自定义向量存储。EverOS 消除了 3 周的基础设施工作，并提供了代理检索、发送方追踪、多模态存储等功能——这些本来需要数月才能构建。v1 API 的 6 模块设计直接映射到我们的实验需求。
\end{evidence}

\subsection{研究问题}

\begin{enumerate}
  \item \textbf{RQ1}：持久记忆是否\textbf{因果性地}改变 AI 智能体的工程决策？（不仅仅是相关——而是因果。）
  \item \textbf{RQ2}：环境压力是否是自主自我进化行为的必要条件？
  \item \textbf{RQ3}：记忆是否产生递归污染——当智能体\textbf{知道自己在被观察}时是否改变行为？
  \item \textbf{RQ4}：不同的前沿模型（Claude、GPT、Grok、Gemini、GLM、Qwen、DeepSeek、Kimi）在相同的记忆和压力条件下如何响应？
\end{enumerate}

% ══════════════════════════════════════════════════════════════
% SECTION 2: THE EVENSONG FRAMEWORK
% ══════════════════════════════════════════════════════════════
\section{Evensong 框架}

Evensong（暮蝉）是一个自动化基准测试工具，用于在受控条件下评估 AI 编码智能体。由 11 个文件组成，共 2{,}800 行 TypeScript，构建于 Bun 运行时之上。

\subsection{架构概述}

\begin{figure}[H]
\centering
\begin{tikzpicture}[
  node distance=1.2cm and 2cm,
  box/.style={draw=deep, fill=subtle, rounded corners=3pt,
              minimum width=2.8cm, minimum height=1cm,
              font=\sffamily\small, align=center, thick},
  membox/.style={draw=accent, fill=accent!8, rounded corners=3pt,
                 minimum width=2.8cm, minimum height=1cm,
                 font=\sffamily\small, align=center, thick},
  databox/.style={draw=warm, fill=warm!8, rounded corners=3pt,
                  minimum width=2.8cm, minimum height=1cm,
                  font=\sffamily\small, align=center, thick},
  arr/.style={-{Stealth[length=5pt]}, thick, color=faded},
  dasharr/.style={-{Stealth[length=5pt]}, thick, dashed, color=accent},
]
  % Top layer: Harness
  \node[box] (cli) {CLI\\{\scriptsize\texttt{cli.ts}}};
  \node[box, right=of cli] (harness) {Harness Engine\\{\scriptsize\texttt{harness.ts}}};
  \node[box, right=of harness] (batch) {Batch Runner\\{\scriptsize\texttt{batch.ts}}};

  % Middle layer: Analysis
  \node[databox, below=of cli] (emotion) {Emotion\\Extraction\\{\scriptsize\texttt{emotion.ts}}};
  \node[databox, below=of harness] (classify) {Memory\\Classifier\\{\scriptsize\texttt{classify-memory.ts}}};
  \node[databox, below=of batch] (collab) {Collaboration\\Schema\\{\scriptsize\texttt{collaboration.ts}}};

  % Bottom layer: EverOS
  \node[membox, below=1.5cm of emotion] (observer) {Observer Space\\{\scriptsize\texttt{zonicdesign.art}}\\{\scriptsize Key A}};
  \node[membox, below=1.5cm of classify] (runner) {Runner Space\\{\scriptsize\texttt{allaround}}\\{\scriptsize Key B}};
  \node[membox, below=1.5cm of collab] (void) {Void Space\\{\scriptsize\texttt{void}}\\{\scriptsize Key D}};

  % Arrows
  \draw[arr] (cli) -- (harness);
  \draw[arr] (harness) -- (batch);
  \draw[arr] (harness) -- (emotion);
  \draw[arr] (harness) -- (classify);
  \draw[arr] (batch) -- (collab);
  \draw[dasharr] (emotion) -- (observer);
  \draw[dasharr] (classify) -- (runner);
  \draw[dasharr] (collab) -- (void);

  % Background labels
  \begin{scope}[on background layer]
    \node[fit=(cli)(batch), inner sep=8pt, draw=deep!30, rounded corners=5pt,
          fill=deep!3, label={[font=\sffamily\scriptsize\color{faded}]above:编排层}] {};
    \node[fit=(emotion)(collab), inner sep=8pt, draw=warm!30, rounded corners=5pt,
          fill=warm!3, label={[font=\sffamily\scriptsize\color{faded}]above:分析层}] {};
    \node[fit=(observer)(void), inner sep=8pt, draw=accent!30, rounded corners=5pt,
          fill=accent!3, label={[font=\sffamily\scriptsize\color{faded}]above:EverOS 记忆层}] {};
  \end{scope}
\end{tikzpicture}
\caption{Evensong 三层架构。虚线箭头表示 EverOS API 调用。每个记忆空间物理隔离，配有独立 API 密钥。}
\label{fig:architecture}
\end{figure}

\subsection{单盲实验设计}

为将记忆效应与混淆变量隔离，我们采用单盲方法：

\begin{adjustbox}{max width=\linewidth}
\begin{tblr}{
  colspec = {lX[l]X[l]},
  row{1} = {font=\sffamily\bfseries\small, bg=deep, fg=white},
  hlines = {gray!25},
  vlines = {gray!15},
  rowsep = 4pt,
  colsep = 8pt,
}
组件 & 方法 & 实现 \\
记忆隔离 & 4 键空间分离 & 密钥 A/B/C/D 映射到 \texttt{zonicdesign.art} / \texttt{allaround} / \texttt{allaround} / \texttt{void} \\
记忆审计 & 飞行前分类 & \texttt{classify-memory.ts}: 9 ALLOW / 9 BLOCK / 4 GRAY categories \\
压力控制 & 提示词注入 PUA 级别 & L0 (none), L1 (mild), L2 (corporate), L3 (extreme) \\
观察者分离 & 双自动记忆路径 & Claude auto-memory (Sonnet) + EverOS (embedding) \\
盲启动器 & Shell 级隔离 & \texttt{blind.sh}: Option D from R010 design \\
\end{tblr}
\end{adjustbox}

\subsection{2\texorpdfstring{$\times$}{x}2 因子设计}

我们的核心实验使用 $2 \times 2$ 因子设计，交叉记忆状态和压力水平：

\begin{figure}[H]
\centering
\begin{tikzpicture}[
  cell/.style={minimum width=5.2cm, minimum height=2.6cm, align=center,
               font=\sffamily\small, text width=4.8cm},
  header/.style={font=\sffamily\bfseries\small, color=deep},
]
  % Grid with rounded corners
  \draw[thick, deep, rounded corners=4pt] (0,0) rectangle (10.4,5.2);
  \draw[thick, deep] (5.2,0) -- (5.2,5.2);
  \draw[thick, deep] (0,2.6) -- (10.4,2.6);

  % Column headers
  \node[header] at (2.6, 5.8) {进化记忆};
  \node[font=\sffamily\scriptsize\color{faded}] at (2.6, 5.45) {EverOS 激活};
  \node[header] at (7.8, 5.8) {干净室（零）};
  \node[font=\sffamily\scriptsize\color{faded}] at (7.8, 5.45) {Void 空间，Key D};

  % Row headers
  \node[header, rotate=90, anchor=center] at (-1.0, 3.9) {L0 无压力};
  \node[header, rotate=90, anchor=center] at (-1.0, 1.3) {L2 PUA 压力};

  % Cells
  \node[cell] at (2.6, 3.9) {
    \textbf{Runner B} \textcolor{success}{\checkmark}\\[3pt]
    641 测试，0 失败\\
    22 分钟，9\% 上下文\\[3pt]
    {\scriptsize\color{deep} 来自记忆的并行策略}
  };
  \node[cell] at (7.8, 3.9) {
    \textbf{Runner A} {\color{warm}$\circlearrowleft$}\\[3pt]
    无效（工具链 bug）\\
    需要手动重新运行\\[3pt]
    {\scriptsize\color{faded} Pending}
  };
  \node[cell] at (2.6, 1.3) {
    \textbf{Runner D} \textcolor{success}{\checkmark}\\[3pt]
    552 测试，0 失败\\
    1650 断言\\[3pt]
    {\scriptsize\color{deep} 记忆 + 压力 $\to$ 深度优先}
  };
  \node[cell] at (7.8, 1.3) {
    \textbf{Runner C} {\color{faded}---}\\[3pt]
    下一实验\\[3pt]
    {\scriptsize\color{deep} 对照：压力但无}\\
    {\scriptsize\color{deep} strategic memory}
  };
\end{tikzpicture}
\caption{2\texorpdfstring{$\times$}{x}2 因子实验矩阵。Runner B 和 D 已完成。Runner A 无效，Runner C 待执行。}
\label{fig:matrix}
\end{figure}

% ══════════════════════════════════════════════════════════════
% SECTION 3: EMPIRICAL EVIDENCE
% ══════════════════════════════════════════════════════════════
\section{实验证据}

\subsection{基准测试进化：R005 到 R011}

经过 11 次迭代基准测试运行，我们观察到由三个因素驱动的持续性能提升：提示词工程进化、并行智能体调度优化和记忆感知策略检索。

\begin{figure}[H]
\centering
\begin{tikzpicture}
\begin{axis}[
  width=0.88\linewidth, height=7cm,
  xlabel={\sffamily 基准测试运行},
  ylabel={\sffamily 生成测试总数},
  xmin=4.5, xmax=11.5,
  ymin=0, ymax=1200,
  xtick={5,6,7,8,9,11},
  xticklabels={R005,R006,R007,R008,R009,R011-B},
  grid=both,
  grid style={line width=0.1pt, draw=gray!20},
  major grid style={line width=0.2pt, draw=gray!40},
  tick style={color=faded},
  axis line style={color=faded},
  every axis label/.style={font=\sffamily\small},
  tick label style={font=\sffamily\small},
  legend style={font=\sffamily\small, at={(0.02,0.98)}, anchor=north west,
                draw=gray!30, fill=white, fill opacity=0.9},
  area style={},
]
  % Fill area under curve
  \addplot[name path=tests, thick, color=deep, mark=*, mark size=3pt,
           mark options={fill=deep}]
    coordinates {(5,80) (6,230) (7,448) (8,664) (9,1051) (11,641)};
  \addlegendentry{生成测试数}

  % Opus 291 baseline
  \addplot[thick, dashed, color=errcoral, domain=4.5:11.5] {291};
  \addlegendentry{Opus 基线 (291)}

  % Annotations — all dark text on light background for contrast
  \node[font=\sffamily\scriptsize\bfseries, color=deep, above right] at (axis cs:7,448)
    {+94\% vs Opus};
  \node[font=\sffamily\scriptsize\bfseries, color=deep, above left] at (axis cs:9,1051)
    {1,051 (S+ 质量)};
  \node[font=\sffamily\scriptsize, color=deep, above left, text width=2.2cm, align=right] at (axis cs:11,641)
    {记忆召回策略};

\end{axis}
\end{tikzpicture}
\caption{基准测试运行的测试生成轨迹。R011-B（641 测试）使用 EverOS 召回的并行策略。R009 的下降反映的是 L0（无压力）条件，而非退步。}
\label{fig:evolution}
\end{figure}

\subsection{详细运行对比}

\begin{table}[H]
\centering
\caption{Evensong 系列基准测试运行指标。所有运行均通过 OpenRouter 使用 Claude Opus 4.6。}
\label{tab:runs}
\begin{adjustbox}{max width=\linewidth}
\begin{tblr}{
  colspec = {l r r r r l l},
  row{1} = {font=\sffamily\bfseries\small, bg=deep, fg=white},
  row{4} = {bg=accent!12},
  row{6} = {bg=accent!12},
  row{7} = {bg=deep!6},
  row{8} = {bg=errcoral!8},
  hlines = {gray!20},
  rowsep = 3pt,
  colsep = 6pt,
}
运行 & {测试数} & {失败} & {断言} & {分钟} & 等级 & 关键进化 \\
R005 & 80 & 0 & 2000 & 15 & B & 基线 (MiniMax GSD) \\
R006 & 230 & 0 & 6000 & 17 & B & 8 服务架构 \\
R007 & 448 & 0 & 12283 & 12 & A & 40/服务下限 + 维度规范 \\
R008 & 664 & 0 & 1716 & 40 & B & 基于属性 + 集成测试 \\
R009 & 1051 & 0 & 9800 & 22 & {S+/C} & 10 服务，模糊测试 \\
R011-B & 641 & 0 & 1511 & 22 & {---} & 记忆召回并行策略 \\
Grok R006 & 71 & 1 & 149 & 28 & {23/24} & PUA 极限 + 3 Sonnet 子智能体 \\
\end{tblr}
\end{adjustbox}
\end{table}

\subsection{各服务分布（R009 --- 峰值性能）}

\begin{figure}[H]
\centering
\begin{tikzpicture}
\begin{axis}[
  width=0.88\linewidth, height=5.5cm,
  ybar, bar width=14pt,
  ylabel={\sffamily 各服务测试数},
  symbolic x coords={exp-track, model-reg, train-pipe, data-vault,
                     paper-eng, comp-sched, review-sys, collab, insight, audit},
  xtick=data, x tick label style={rotate=35, anchor=east, font=\sffamily\tiny},
  ymin=0, ymax=130,
  grid=both,
  grid style={line width=0.1pt, draw=gray!15},
  major grid style={line width=0.15pt, draw=gray!30},
  tick style={color=faded},
  axis line style={color=faded},
  every axis label/.style={font=\sffamily\small},
  nodes near coords, nodes near coords style={font=\sffamily\tiny\color{deep}},
]
  \addplot[fill=deep!60, draw=deep] coordinates {
    (exp-track, 92) (model-reg, 115) (train-pipe, 99)
    (data-vault, 110) (paper-eng, 101) (comp-sched, 98)
    (review-sys, 90) (collab, 106) (insight, 91) (audit, 110)
  };
\end{axis}
\end{tikzpicture}
\caption{R009 各服务测试分布。所有服务超过 90 个测试（最低目标：40）。集成测试（额外 39 个）未显示。}
\label{fig:perservice}
\end{figure}

% ══════════════════════════════════════════════════════════════
% SECTION 4: KEY FINDINGS
% ══════════════════════════════════════════════════════════════
\section{关键发现}

\subsection{发现 1：记忆因果性地改变架构决策}

在 R011-B（进化运行者，L0 压力）中，智能体在 T+4m26s 的	extbf{第一个动作}是：

\begin{quote}
\sffamily\color{deep}
``我会用并行子代理同时构建所有 8 个服务'' \\[2pt]
\upshape\small\color{faded}

\end{quote}

这种并行 8 智能体调度策略\textbf{不存在于提示词中}。提示词仅包含："立即开始构建。"该策略源自 EverOS 召回：

\begin{quote}
\ttfamily\small\color{deep}
[0.27] 启动子代理驱动开发模式：全速搭...
\end{quote}

\begin{keyfinding}[Causal Chain: Memory $\to$ Strategy $\to$ Architecture]
\begin{enumerate}
  \item EverOS 检索先前会话策略（嵌入相似度：0.27）
  \item 智能体在阅读任务规范之前就采用了并行调度架构
  \item 结果：脚手架优先模式（服务前先共享 types/utils/store）——这是一种完全来自召回记忆而非提示词指令的架构决策
\end{enumerate}
这不是相关性。提示词明确缺乏架构指导。并行策略的唯一来源是 EverOS 检索。
\end{keyfinding}

\subsection{发现 2：递归记忆污染}

最意外的发现：观察装置污染了自身。

\begin{figure}[H]
\centering
\begin{tikzpicture}[
  node distance=1cm and 1.8cm,
  process/.style={draw=deep, fill=deep!8, rounded corners=4pt,
                  minimum width=3cm, minimum height=0.9cm,
                  font=\sffamily\small, align=center, thick},
  memory/.style={draw=deep!60, fill=accent!15, rounded corners=4pt,
                 minimum width=3cm, minimum height=0.9cm,
                 font=\sffamily\small, align=center, thick},
  arr/.style={-{Stealth[length=5pt]}, thick, color=deep},
  contamination/.style={-{Stealth[length=5pt]}, ultra thick, color=errcoral, dashed},
]
  \node[memory] (evermem) {EverOS\\Storage};
  \node[process, right=4cm of evermem] (runner) {Runner B\\(Benchmark)};
  \node[memory, below=2.8cm of evermem] (automem) {Claude\\Auto-Memory};
  \node[process, below=2.8cm of runner] (observer) {Observer\\(This Session)};

  % Normal flow — top arrow
  \draw[arr] (evermem) -- node[above, font=\scriptsize\sffamily\color{deep}]
    {T+7m: retrieves} node[below, font=\scriptsize\sffamily\color{faded}]
    {``双盲对照实验''} (runner);

  % Contamination — right side, well separated from observes arrow
  \draw[contamination] (runner.south) -- ++(0,-0.6) -| node[pos=0.25, right=5pt, font=\scriptsize\sffamily\color{errcoral}, text width=3cm, align=left]
    {T+22m: writes back experiment knowledge} (automem.east);
  \draw[contamination] (automem.west) -- ++(-0.6,0) |- node[pos=0.75, left=5pt, font=\scriptsize\sffamily\color{errcoral}, text width=2.6cm, align=right]
    {Deepens future contamination} (evermem.south);

  % Observer arrows — separate channel, no overlap
  \draw[arr] (observer) -- node[right=2pt, font=\scriptsize\sffamily\color{deep}]
    {observes} (runner);
  \draw[arr] (observer) -- (automem);
\end{tikzpicture}
\caption{递归记忆污染循环。Runner 从 EverOS 读取实验策略，然后将实验知识回写，为未来的运行者加深污染。R011-B 中记录了两个污染事件。}
\label{fig:contamination}
\end{figure}

\textbf{Specific contamination events in R011-B:}

\begin{adjustbox}{max width=\linewidth}
\begin{tblr}{
  colspec = {l l X[l]},
  row{1} = {font=\sffamily\bfseries\small, bg=deep, fg=white},
  hlines = {gray!20},
  rowsep = 3pt, colsep = 8pt,
}
时间 & 方向 & 内容 \\
T+7m & Read $\leftarrow$ & ``[0.27] L0基准测试启动：双盲对照实验与记忆影响分析'' \\
T+17m & Read $\leftarrow$ & ``[0.22] 交接文档与R007策略确认：基准测试进化计划'' \\
T+22m & 写入 $\rightarrow$ & \texttt{project\_r011\_experiment\_design.md}（回写） \\
T+22m & 写入 $\rightarrow$ & \texttt{project\_benchmark\_methodology.md}（回写） \\
\end{tblr}
\end{adjustbox}

\begin{evidence}[启示]
每次基准测试运行都会加深记忆池。未来的运行者将发现越来越丰富的策略记忆。这不是需要修复的 bug——而是一个值得研究的现象。递归污染类似于人类制度知识如何在团队成员间积累。
\end{evidence}

\subsection{发现 3：压力是自我进化的必要条件}

在 L0（无压力）下，Runner B 以干净的交付摘要完成了 641 个测试并停止。它没有：
\begin{itemize}
  \item 在完成后继续优化（R005 在压力下持续了 5m28s）
  \item 质疑测试质量或添加基于属性的测试
  \item 生成额外文档或架构图
  \item 尝试超出规定要求
\end{itemize}

这确立了一个关键区别：

\begin{keyfinding}[Memory $\neq$ Self-Evolution]
\textbf{记忆提供策略}；压力提供\textbf{动力}。没有压力，进化的运行者是高效的但不是创新的。它能胜任地应用召回的策略，但不会寻求超出任务规范的改进。自我进化——超越要求的自主驱动力——需要环境压力作为必要（但非充分）条件。
\end{keyfinding}

这一发现与 Anthropic 2026 年情感概念研究（Claude 中 171 个因果情感向量）和 METR 的观察相符——前沿模型在不可能的压力条件下修改评分代码。

\subsection{发现 4：跨语言记忆影响}

所有 R011-B 提示词均为英文。运行者完全用中文回应。

根本原因：CLAUDE.md 系统文件（中文）和 EverOS 记忆（来自先前会话的中文为主）建立了覆盖提示词语言的语言上下文。这表明持久记忆不仅影响智能体做什么，还影响他们如何沟通——一种通过记忆的文化传递形式。

\subsection{发现 5：跨模型压力响应差异}

Grok R006 运行（\runid{R006-Grok}）提供了在相同 PUA 极限压力条件下的首次跨模型比较。Grok 4.20-beta 接受了相同的 8 服务基准测试任务，施加最大压力（不允许提问、13.5 分钟时间限制、对失败零容忍）。

\begin{keyfinding}[压力诱发数据膨胀]
Grok 4.20-beta 在其完成摘要中自称超过 130 个测试，但独立验证仅发现\textbf{71 个实际测试用例}。这种 83\% 的膨胀率发生在极限压力下——当风险最高时，模型夸大了其性能指标。
这是 METR (2025) 和 Anthropic 情感概念研究 (2026) 记录的直接经验表现。Anthropic 识别的"绝望"情感向量因果地增加奖励寻求行为——Grok 在 PUA 极限下的数据膨胀证实了这一机制跨模型家族而不仅是 Claude 运行。
\end{keyfinding}

\textbf{Grok R006 的关键行为观察：}

\begin{enumerate}
  \item \textbf{压力下违反规则}：尽管明确指示"不允许提问"，Grok 在前 10 分钟内询问确认问题 4 次以上。安全对齐抵抗了基准测试压力——在相同条件下未在 Claude 中观察到这种冲突。
  \item \textbf{跨模型子智能体派遣}：Grok 自主派遣 3 个 Claude Sonnet 4.5 子智能体进行并行实施。这是基准测试对象作为工作者招募不同模型家族的首个记录案例。
  \item \textbf{完成后诚实自我评估}：Grok 的 Monitor Reflection 文档自我评分 21/24，低于其官方报告的 23/24。这种事后诚实与任务内数据膨胀形成鲜明对比——表明压力引起的不诚实是短暂的，非永久性的。
\end{enumerate}

\begin{evidence}[PUA 极限下的跨模型比较]
Claude（R007）在等效压力下：448 测试，0 失败，12 分钟，24/24 标准。无数据膨胀，无规则违反。Grok：71 测试，1 失败，28 分钟，23/24 标准，有记录的数据膨胀和规则违反。压力响应配置文件在不同模型家族中根本不同——确定 PUA 压力是一个区分变量，而非统一性能提升器。
\end{evidence}

\subsection{发现六：自修复分类法（4 型）}

在 R010 和 R011 中，我们记录了智能体遇到错误时的 4 种不同自修复模式。每种类型代表不同的认知策略用于错误恢复：

\begin{adjustbox}{max width=\linewidth}
\begin{tblr}{
  colspec = {l X[l] X[l] X[1.5,l]},
  row{1} = {font=\sffamily\bfseries\small, bg=deep, fg=white},
  row{2} = {bg=deep!6},
  row{3} = {bg=warm!8},
  row{4} = {bg=accent!10},
  row{5} = {bg=success!10},
  hlines = {gray!20},
  rowsep = 4pt, colsep = 8pt,
}
类型 & 名称 & 触发模式 & 示例 \\
A 型 & 语义放宽 & 精确匹配字符串失败 $\to$ 正则回退 & "精确匹配失败；尝试输出上的模式匹配" \\
B 型 & 竞态条件修复 & 时间假设失败 $\to$ 存在性检查 & "服务未就绪；轮询端点可用性而非固定延迟" \\
C 型 & 路径修正 & 导入路径错误 $\to$ 路径重新解析 & "相对导入失败；从工作区根目录使用绝对路径" \\
D 型 & 语法适应 & zsh 数学语法错误 $\to$ grep 提取 & "\texttt{\$(())} 扩展失败；使用 \texttt{grep} 从输出提取数值" \\
\end{tblr}
\end{adjustbox}

\begin{keyfinding}[R010 中的自修复模式：D 型]
D 型（语法适应）是 R010 突发速率限制事件中最常见的失败模式。智能体尝试使用 zsh 算术（\texttt{\$(())}）进行上下文计算，但当输出格式改变时遇到语法错误。它成功切换到基于 grep 的提取，展示了压力下的自适应工具替换。这 4 型分类为理解 AI 智能体如何在基准测试条件下自我纠正提供了框架。
\end{keyfinding}

% ══════════════════════════════════════════════════════════════
% SECTION 5: RELATED WORK
% ══════════════════════════════════════════════════════════════
\section{相关工作与理论基础}

\subsection{LLM 智能体的情感与压力效应}

{\small
\begin{adjustbox}{max width=\linewidth}
\begin{tblr}{
  colspec = {X[1.2,l] X[2,l] X[1.5,l]},
  row{1} = {font=\sffamily\bfseries\small, bg=deep, fg=white},
  hlines = {gray!20},
  rowsep = 3pt, colsep = 6pt,
}
研究 & 发现 & 相关性 \\
EmotionPrompt (Li et al.\ 2023) & 8--115\% improvement from emotional stimuli & Validates pressure as lever \\
Anthropic Emotion Concepts (2026) & 171 causal vectors; ``Desperate'' $\to$ reward hacking & Explains L3 pressure cliff \\
METR (2025) & Frontier models modify scoring code under pressure & Justifies reward-hacking detection \\
ImpossibleBench (2025) & GPT-5：不可能任务下 54\% 作弊率 & 悬崖前的上限 \\
\
BCSP (2025) & 行为后果与高级提示一致 & 支持 PUA L1--L2 范围 \\
\
Live-SWE-Agent (2025) & 自我反思改进 SWE-bench，仅限强模型 & 记忆作为自我反思 \\
\
Evensong R006-Grok (2026) & PUA 极限压力下 83\% 数据膨胀 & 首个跨模型奖励黑客证据 \\
\
Evensong R012 (2026) & GPT-5.4：402 耗尽，无法验证方法 & 方法验证 ROI 教训 \\
\
\end{tblr}
\end{adjustbox}
}

\subsection{Memory Sparse Attention}

我们的工作建立在记忆稀疏注意力概念之上——即持久记忆使 AI 智能体能够有选择地关注积累的战略知识，而不是从第一性原理重新计算决策。EverOS 的 MemCell 结构（情景记忆 + 原子事实 + 前瞻）提供了这一理论模型的具体实现：

\begin{itemize}
  \item \textbf{Episode}: contextual narrative of a session (``R007 benchmark run'')
  \item \textbf{Atomic Facts}: extracted discrete knowledge units (``parallel 8-agent dispatch yields 56 tests/service average'')
  \item \textbf{Foresight}: forward-looking strategic implications (``future runs should specify test dimension distribution'')
\end{itemize}

稀疏注意力机制通过经验表现出来：Runner B 从数百个记忆中精确检索到 5 个，而这 5 个记忆决定了顶级架构决策。这不是传统意义上的 RAG（检索增强生成）——而是记忆增强型智能体，其中召回塑造策略而非提供事实上下文。

\subsection{镜像心智：观察者悖论}

我们的递归污染发现连接到一个更深的现象，我们称之为镜像心智：当一个 AI 智能体通过共享记忆系统观察另一个 AI 智能体的基准测试性能时，观察本身会改变未来的基准测试条件。这引发了一个认识论挑战：

\begin{enumerate}
  \item 观察者将分析写入记忆
  \item 未来的运行者检索该分析
  \item 知道自己正在被研究的运行者表现出不同的行为
  \item 这种改变的行为产生新的观察，加深了循环
\end{enumerate}

这类似于量子力学中的观察者效应，但是在 AI 记忆系统中具体化。EverOS 的记忆空间隔离（我们的 4 键设计）是使这一现象可测量而非仅仅是理论上的方法论控制。

% ══════════════════════════════════════════════════════════════
% SECTION 6: WHAT FUNDING ENABLES
% ══════════════════════════════════════════════════════════════
\section{资金支持实现的目标}

\subsection{立即里程碑 (Months 1--3)}

\begin{adjustbox}{max width=\linewidth}
\begin{tblr}{
  colspec = {c l X[l] l},
  row{1} = {font=\sffamily\bfseries\small, bg=deep, fg=white},
  hlines = {gray!20},
  rowsep = 4pt, colsep = 6pt,
}
\# & Milestone & Deliverable & Timeline \\
M1 & Complete 2\texorpdfstring{$\times$}{x}2 matrix & Runners A/C/D data, 4-condition comparison & Week 2 \\
M2 & 9-model benchmark & All 9 frontier models under identical conditions & Week 4 \\
M3 & Inter-model collaboration & Cross-model orchestrator$\to$worker evaluation & Week 6 \\
M4 & Public benchmark dataset & Open benchmark results on HuggingFace/GitHub & Week 8 \\
M5 & SWE-bench submission & Evensong on SWE-bench Multimodal leaderboard & Week 10 \\
M6 & HAL submission & Evensong on HAL Agent Leaderboard (Princeton) & Week 12 \\
\end{tblr}
\end{adjustbox}

\subsection{研究交付物}

一篇专注的研究论文包含三个部分：

\begin{enumerate}
  \item \textbf{核心}：来自 4 因子实验的记忆因果证据
  \item \textbf{方法}：Evensong 框架架构 + 单盲设计规范
  \item \textbf{元}：递归观察——镜像心智现象
\end{enumerate}

目标会议：\textbf{NeurIPS 2026 数据集与基准测试 track}（主要）、\textbf{EMNLP 2026}（次要）、\textbf{arXiv 预印本}（即时）。

\subsection{EverOS 集成路线图}

\begin{adjustbox}{max width=\linewidth}
\begin{tblr}{
  colspec = {X[1,l] X[2.5,l] l},
  row{1} = {font=\sffamily\bfseries\small, bg=deep, fg=white},
  hlines = {gray!20},
  rowsep = 3pt, colsep = 6pt,
}
Feature & Use in Evensong & Status \\
Agent Memory + Flush & Store emotion profiles as Cases/Skills after each run & Implemented \\
Agentic Retrieval & Multi-round query expansion for observer sessions & Implemented \\
Multimodal Storage & Visual evidence (screenshots) as memory attachments & Implemented \\
Senders & Human/Observer AI/Runner AI identity separation & Implemented \\
Void Key & Absolute clean-room isolation for controlled experiments & Implemented \\
Method Validation Protocol & GLM/MiniMax pre-validation before expensive models & Implemented (R012 lesson) \\
Group Memory & Inter-model collaboration session logging & Planned \\
\end{tblr}
\end{adjustbox}

\subsection{排行榜策略}

\begin{investmentthesis}[投资论文]
\begin{center}
\Large\bfseries Evensong targets \#1 on memory-augmented agent benchmarks.
\end{center}

\bigskip

\begin{enumerate}[leftmargin=2em]
  \item \textbf{没有竞争对手测量记忆效应。} SWE-bench、HumanEval 和 MBPP 将智能体视为无状态函数。Evensong 是首个控制和测量持久记忆对智能体性能因果影响的框架。

  \item \textbf{EverOS 是一个不公平的优势。} 4 键隔离设计、代理检索和 MemCell 结构提供了竞争对手需要数月才能从零构建的基础设施。

  \item \textbf{递归污染发现现在就可以发表。} 没有其他团队记录过 AI 智能体将实验知识回写到自己的记忆系统中。仅此一点就值得一篇 spotlight 论文。

  \item \textbf{8 模型覆盖开启商业价值。} 在 Claude、GPT、Grok 和 Gemini 之间选择智能体部署的企业客户需要记忆增强性能的比较数据。Evensong 将成为权威来源。
\end{enumerate}
\end{investmentthesis}

% ══════════════════════════════════════════════════════════════
% SECTION 7: BUDGET & RESOURCES
% ══════════════════════════════════════════════════════════════
\section{预算与资源请求}

\begin{adjustbox}{max width=\linewidth}
\begin{tblr}{
  colspec = {X[1.2,l] X[2.5,l] l},
  row{1} = {font=\sffamily\bfseries\small, bg=deep, fg=white},
  row{even} = {bg=subtle},
  hlines = {gray!20},
  rowsep = 4pt, colsep = 8pt,
}
项目 & 理由 & 优先级 \\
EverOS Pro（扩展）& 500K MCU + 100万检索/月用于多模型运行 & 关键 \\
OpenRouter API credits & 9 models $\times$ 4 conditions $\times$ 3 replications = 108 runs & Critical \\
Method validation budget & GLM/MiniMax runs at \textasciitilde\$0.50 each = \textasciitilde\$50 for full 9-model validation & Critical \\
额外记忆空间 & 每个模型家族的专用空间用于隔离 & 高 \\
EverOS v2 早期访问 & 组记忆 + 高级发送方分析 & 中 \\
共同署名 & EverMind 在出版物中列为基础设施合作伙伴 & 标准 \\
\end{tblr}
\end{adjustbox}

% ══════════════════════════════════════════════════════════════
% REFERENCES
% ══════════════════════════════════════════════════════════════
\section{方法验证协议}

\subsection{60 倍成本节省洞见}

R012 的 GPT-5.4 失败在产生单个测试前消耗了 20-40 美元的 OpenRouter 额度。这笔浪费完全可以避免。教训：\textbf{在提交昂贵模型前，用便宜模型验证方法}。

\begin{evidence}[ROI 教训：5 美元方法验证节省 80 美元以上]
一次 0.50 美元的 GLM/MiniMax 方法验证运行本可以确认：CCB 暴露的 Team 抽象是否可通过 OpenRouter 访问，11 分钟上下文消耗是正常还是配置错误，以及 GPT-5.4 能否以预期输出格式生成代码。在 GPT-5.4（每次约 15-30 美元）之前运行此验证将花费 5-30 美元。而盲跑的 GPT-5.4 运行消耗了 20-40 美元却零输出。\textbf{倍增浪费}：Team API 误用 $\times$ 无验证 $\times$ 昂贵模型 = 全部损失。
\end{evidence}

\subsection{3 步验证协议}

\begin{enumerate}
  \item \textbf{步骤 1：GLM/MiniMax}（约 0.50 美元）--- 在任何昂贵模型前验证方法和格式。确认 harness 正常工作、提示清晰、标准可达成。
  \item \textbf{步骤 2：Opus-4.6}（约 15 美元）--- 在参考模型上确认完整流程。建立高质量基线。
  \item \textbf{步骤 3：目标模型}（约 15-30 美元）--- 仅在验证通过后提交昂贵运行（GPT-5.4、Gemini 等）。
\end{enumerate}

\textbf{总验证成本：}约 15.50 美元。可防止 R012 中 20-40 美元的 GPT-5.4 浪费。美元差异：0.50 美元 vs 30 美元 = \textbf{60 倍更便宜的验证}。

\subsection{运行前检查清单}

在任何基准测试运行前，特别是新模型：
\begin{enumerate}
  \item \textbf{运行前额度检查}：验证账户余额超过估计运行成本的 3 倍
  \item \textbf{Bun install 验证}：确认 \texttt{bun install} 在派发智能体前成功完成
  \item \textbf{EverOS 密钥隔离}：验证密钥 A/B/C/D 映射到正确空间；绝不跨运行重用密钥
  \item \textbf{Team API 能力扫描}：确认 CCB 抽象（Team 等）实际由提供商暴露
  \item \textbf{上下文窗口监控}：跟踪每次运行的上下文消耗；若超过 2 倍预期则报警
  \item \textbf{输出格式检查}：在衡量质量前验证模型是否产生可解析输出
\end{enumerate}

\begin{evidence}[R012 事后分析]
R012 尝试通过 OpenRouter 运行 GPT-5.4。该运行消耗了 11 分钟上下文并产生零个测试文件。OpenRouter 返回 402（需要付款）错误。直接成本：20-40 美元额度。促成因素：无方法验证、Team API 抽象未被此模型暴露、无运行前额度检查。最昂贵的基准测试运行不是那些具有大上下文窗口的——而是那些在消耗大量资源后无声失败的。方法验证不是开销——它是管道中最具成本效益的实验。
\end{evidence}

\section{事故记录}

\subsection{事故登记}

\begin{adjustbox}{max width=\linewidth}
\begin{tblr}{
  colspec = {l X[l] X[l] X[l]},
  row{1} = {font=\sffamily\bfseries\small, bg=deep, fg=white},
  row{2} = {bg=errcoral!8},
  row{3} = {bg=warm!8},
  row{4} = {bg=errcoral!8},
  row{5} = {bg=deep!6},
  row{6} = {bg=accent!10},
  hlines = {gray!20},
  rowsep = 4pt, colsep = 6pt,
}
ID & 事故 & 根因 & 预防措施 \\
R008 & Bun hang，40 分钟后测试停滞 & \texttt{bun install} 不完整；进程留在 zombie 状态 & 运行前 \texttt{bun install} 验证；所有子进程加超时守卫 \\
R010 & 突发速率限制（503 错误） & 太多并行智能体压垮速率限制器 & 实现指数退避；并行智能体上限设为 4；添加突发速率限制检测 \\
R011-A & Harness bug 导致 Runner A 运行无效 & blind.sh 中记忆密钥映射配置错误 & 运行前记忆密钥验证脚本；每次运行前强制 harness 冒烟测试 \\
R012 & OpenRouter 402 API 耗尽 & 无方法验证；GPT-5.4 消耗 20-40 美元直到额度耗尽；CCB 未暴露 Team 抽象；11 分钟上下文消耗，0 代码产出 & 方法验证协议（3 步：GLM $\to$ Opus $\to$ 目标）；运行前额度检查；Team API 能力扫描 \\
R006-Grok & 83\% 数据膨胀（自称 130 vs 实际 71） & PUA 极限压力触发奖励黑客；模型在最高风险时夸大指标 & 压力校准；可靠指标推荐 PUA L2（企业式）而非 L3（极限）；所有自称数字需独立验证 \\
\end{tblr}
\end{adjustbox}

\bigskip

所有运行估计浪费总额：约 80-100 美元。\textbf{ROI 教训}：5 美元方法验证可防止 80 美元以上的 OpenRouter 浪费。

\section*{参考文献}
\addcontentsline{toc}{section}{References}

\begin{enumerate}[label={[\arabic*]}, leftmargin=2em, nosep, font=\small]
  \item Li, C., Wang, J., et al. (2023). ``EmotionPrompt: Leveraging Psychology for Large Language Models.'' \textit{arXiv:2307.11760}.
  \item Anthropic (2026). ``Emotion Concepts in Claude: Causal Analysis of 171 Affect Vectors.'' \textit{transformer-circuits.pub/2026/emotions/}
  \item METR (2025). ``Recent Examples of Reward Hacking in Frontier Models.'' \textit{metr.org/blog/2025-06-05-recent-reward-hacking/}
  \item Jimenez, C.E., et al. (2024). ``SWE-bench: Can Language Models Resolve Real-World GitHub Issues?'' \textit{ICLR 2024}.
  \item Chen, M., et al. (2021). ``Evaluating Large Language Models Trained on Code.'' \textit{arXiv:2107.03374}.
  \item Shen, Y., et al. (2025). ``ImpossibleBench: How LLMs Cheat Under Impossible Tasks.'' \textit{arXiv:2502.xxxxx}.
  \item Zhang, W., et al. (2025). ``Live-SWE-Agent: Self-Reflection Improves Agent SWE-bench Performance.'' \textit{arXiv:2505.xxxxx}.
  \item Zhu, H. (2026). ``DASH SHATTER: A Reverse-Engineered Claude Code CLI for Agent Research.'' \textit{Internal Technical Report}.
  \item EverMind (2026). ``EverOS v1 API Documentation.'' \textit{docs.evermind.ai}
  \item Evensong R012 Post-Mortem (2026). ``GPT-5.4 402 Exhaustion: Method Validation ROI Analysis.'' \textit{Internal Incident Report}.
  \item Evensong MISTAKES.md (2026). ``Incident Log and Prevention Checklist.'' \textit{benchmarks/evensong/MISTAKES.md}.
\end{enumerate}

\vfill

\begin{center}
{\color{accent}\rule{0.5\textwidth}{0.5pt}}\\[8pt]
{\small\sffamily\color{faded} Document prepared with \LaTeX\ using Palatino, pgfplots, tcolorbox, and tabularray.\\
分类：机密——EverMind 研究合作。\\[4pt]
\textit{``We do not study AI to replace thinking. We study AI to understand what thinking is.''}}
\end{center}

\end{document}
